% Part: first-order-logic
% Chapter: models-theories
% Section: expressing-props-of-structures

\documentclass[../../../include/open-logic-section]{subfiles}

\begin{document}

\olfileid{fol}{mat}{exs}

\olsection{التعبير عن خصائص ال\printtoken{P}{structure}}

\begin{explain}
كثيرًا ما يكون من المفيد والمهم أن نعبر عن شروط على الدوال والعلاقات،
أو، بوجه أعم، عن كون الدوال والعلاقات في بنية تحقق تلك الشروط. فنود،
مثلًا، أن تكون لدينا سبل لتمييز ال!!{structure}s التي «تجسد» ما نريد
أن «تعنيه» ال!!{predicate}s في لغة ما من ال!!{structure}s التي لا
تفعل ذلك. ولنا، بالطبع، مطلق الحرية في تحديد ال!!{structure}s التي
«نقصدها»؛ فيمكننا، مثلًا، أن نشترط أن يكون تأويل
ال!!{predicate}~$\le$ ترتيبًا، أو أن نقصر اهتمامنا على تأويلات~$\Lang
L$ التي يتألف مجالها من مجموعات ويؤول فيها $\Obj \in$ بعلاقة «هو
!!a{element} في». ولكن هل نستطيع فعل ذلك ب!!{sentence}s اللغة؟ أي:
ما الشروط على !!a{structure}~$\Struct M$ التي نستطيع التعبير عنها
ب!!a{sentence} (أو ربما بمجموعة من ال!!{sentence}s) في لغة~$\Struct
M$؟ ثمة شروط لن نستطيع التعبير عنها. فلا توجد، مثلًا، جملة من~$\Lang
L_A$ لا تصدق في !!a{structure}~$\Struct M$ إلا إذا كان
$\Domain M = \Nat$. ولا نستطيع التعبير عن أن «المجال لا يحتوي إلا
على الأعداد الطبيعية». ولكن ربما استطعنا التعبير عن «خصائص بنيوية»
لل!!{structure}s. فأي خصائص لل!!{structure}s نستطيع التعبير عنها
ب!!{sentence}s؟ أو، بصياغة أخرى، أي مجموعات من ال!!{structure}s
نستطيع وصفها بأنها تلك التي تجعل !!a{sentence} (أو مجموعة من
ال!!{sentence}s) صادقة؟
\end{explain}

\begin{defn}[نموذج مجموعة]
لتكن $\Gamma$ مجموعة من ال!!{sentence}s في لغة~$\Lang L$. نقول إن
!!a{structure}~$\Struct M$ \emph{نموذج لـ}~$\Gamma$ إذا كان
$\Sat{M}{!A}$ لكل $!A \in \Gamma$.
\end{defn}

\begin{ex}
تصدق الجملة $\lforall[x][x \le x]$ في~$\Struct M$ إذا وفقط إذا كانت
$\Assign{\le}{M}$ علاقة انعكاسية. وتصدق الجملة
$\lforall[x][\lforall[y][((x \le y \land y \le x) \lif x = y)]]$ في
~$\Struct M$ إذا وفقط إذا كانت $\Assign{\le}{M}$ مضادة للتناظر.
وتصدق الجملة $\lforall[x][\lforall[y][\lforall[z][((x \le y \land y \le z)
      \lif x \le z)]]]$ في~$\Struct M$ إذا وفقط إذا كانت
$\Assign{\le}{M}$ متعدية. ومن ثم فإن نماذج
\begin{align*}
\{\quad &\lforall[x][x \le x], \\
   & \lforall[x][\lforall[y][((x \le y \land y \le
    x) \lif x = y)]], \\
   &\lforall[x][\lforall[y][\lforall[z][((x \le y
      \land y \le z) \lif x \le z)]]] \quad \}
\end{align*}
هي بالضبط البنى التي تكون فيها~$\Assign{\le}{M}$ انعكاسية ومضادة
للتناظر ومتعدية، أي ترتيبًا جزئيًا. ولذلك يمكننا اتخاذ هذه الجمل
بديهيات \emph{لنظرية الرتبة الأولى للترتيبات الجزئية}.
\end{ex}

\end{document}
